\documentclass[11pt,letterpaper]{article}

\usepackage[margin=0.85in]{geometry}
\usepackage[T1]{fontenc}
\usepackage[utf8]{inputenc}
\usepackage{lmodern}
\usepackage{amsmath,amssymb}
\usepackage{booktabs}
\usepackage{tabularx}
\usepackage{array}
\usepackage{color}
\usepackage{graphicx}
\graphicspath{{../shared/figures/}}
\usepackage[numbers,sort&compress]{natbib}
\usepackage[colorlinks=true,linkcolor=blue,citecolor=blue,urlcolor=blue]{hyperref}

\definecolor{barblue}{rgb}{0.17,0.41,0.66}
\definecolor{bargray}{rgb}{0.55,0.60,0.64}
\definecolor{barteal}{rgb}{0.13,0.55,0.50}
\definecolor{baramber}{rgb}{0.80,0.47,0.09}
\definecolor{barred}{rgb}{0.75,0.18,0.18}

\title{\textbf{How Are Engineering Papers Adopting Lean?}\\
\large An Empirical Longitudinal Measurement of arXiv Preprints, Authoring Paradigms, Recall Calibration, and Semantic Fidelity}

\author{
  Scott Moeller \and Bhaskar Krishnamachari\thanks{Autonomous Networks Research Group (ANRG), University of Southern California. Working draft; pipeline code, structured JSON artifacts, and audit scripts are maintained in \texttt{survey-lean-proofs-eng-arxiv}.}
}
\date{November 23, 2026 \quad (Pilot Report v0.5 --- 9,837-Paper Longitudinal Curve, 23-Repo Layer-A Audit \& 14-Repo Sandboxed Kernel Verification)}

\renewcommand{\topfraction}{0.88}
\renewcommand{\textfraction}{0.10}

\begin{document}
\sloppy
\maketitle

\begin{abstract}
The rapid rise of LLM-assisted autoformalization in 2026 has brought the Lean~4 theorem prover into applied engineering, optimization, and systems research---domains far removed from pure mathematics and traditional formal-methods venues. Yet postmortems of bulk AI-generated Lean releases show that a green Lean build alone does not guarantee semantic fidelity between a paper's prose claims and its formal code, nor does it prevent proof bloat, scope downgrading, or silent hypothesis strengthening. We present a reproducible, token-frugal empirical longitudinal measurement and semantic-fidelity audit of \textbf{9,837} unique arXiv preprints across three matched 19-day windows---\textbf{T0} (Nov~1--19, 2025; $n=2{,}435$), \textbf{T1} (May~1--19, 2026; $n=3{,}298$), and \textbf{T2} (Nov~1--19, 2026; $n=4{,}104$)---spanning \textbf{eight} engineering, optimization, systems, control, information theory, signal processing, security, and robotics categories (\texttt{math.OC}, \texttt{eess.SY}, \texttt{eess.SP}, \texttt{cs.IT}, \texttt{cs.DC}, \texttt{cs.NI}, \texttt{cs.CR}, \texttt{cs.RO}). Combining bias-free OAI-PMH \texttt{arXivRaw} ingestion (omitting \texttt{until} to rescue 31.8--35.4\% later-revised historical preprints), polite full-text HTML5/PDF screening (99.68\% coverage), strictly static LaTeX and GitHub/GitLab/Zenodo artifact profiling, and parallel lightweight-agent classification reduces LLM inspection volume by 98.7\% while achieving $3.4\times$ higher recall than metadata-only search. To referee-proof Stage-2 recall, we audited \textbf{200} random GitHub-linked T2 negatives (333 unique repositories via recursive GitHub API tree queries) plus 100 non-GitHub negatives, finding \textbf{0/200} missed Lean repositories (95\% Clopper--Pearson upper bound: \textbf{1.49\%}).

Across the 12-month longitudinal curve, contemporaneous (\texttt{v1}) Lean adoption surged from \textbf{0/2,435 (0.00\%, 95\% UB 0.12\%)} in Nov~2025 to \textbf{5/3,298 (0.15\%)} in May~2026 to \textbf{39/4,104 (0.95\%)} in Nov~2026---\textbf{a $6.3\times$ jump in six months} with non-overlapping 95\% confidence intervals---while deterministic Stage-2 hits for all \emph{other} formal verification tools combined (Coq/Rocq, Isabelle, TLA+, Tamarin, ProVerif, SMT) remained flat-to-declining (\textbf{1.52\% $\rightarrow$ 1.33\% $\rightarrow$ 1.17\%}), with Lean~4 candidates overtaking all other formal tools combined in Nov~2026 (55 vs.\ 48). Direct \texttt{v1} inspection further uncovers \emph{retroactive post-hoc autoformalization}: three historical preprints (\texttt{2511.15517}, \texttt{2605.09552}, \texttt{2605.19531}) had zero Lean mentions at initial posting but added 374.6k~LoC of Lean~4 in June--Nov~2026 \texttt{v2} revisions. Over the same six months (T1$\rightarrow$T2 contemporaneous \texttt{v1}), median paper-owned Lean codebase size jumped \textbf{38.7$\times$} (749 $\rightarrow$ 29,005~LoC), explicit AI disclosure rose from \textbf{40\%} to \textbf{76.9\%}, and \texttt{Tier3\_Gold} prevalence rose from \textbf{0\%} to \textbf{48.7\%}. In T2, adoption bifurcates sharply: \textbf{\texttt{math.OC} leads at 2.62\% (19/724, 17 \texttt{engineering\_math\_proof})}, followed by \texttt{cs.DC} (1.79\%), \texttt{cs.IT} (1.78\%), \texttt{cs.CR} (1.11\%), and \texttt{eess.SY} (0.41\%), whereas \texttt{cs.NI} (0/559 across 3 slices) and \texttt{cs.RO} (0/2,431 across 3 slices) remain at 0.00\% throughout. Finally, an expanded \textbf{Phase~2 adversarial Layer-A semantic fidelity audit} of \textbf{23} locally inspected \texttt{engineering\_math\_proof} repositories (20 in T2 + 3 in T1) shows both temporal maturation---May~2026 \texttt{v1} repos (473--749~LoC) were \texttt{partial} with moderate Mode~\#3 custom abstractions, whereas \textbf{20/20} Oct--Nov~2026 author-owned repos exhibit zero moderate/severe Mode~\#1 drift or Mode~\#3 smuggling---and a sharp natural control contrast against the sole \texttt{material\_undisclosed\_drift\_or\_risk} failure: bulk automated pipeline \texttt{2611.26272} (\emph{Paper2Lean}/\texttt{OptPaper-20}). To ground-truth static verification tiers dynamically, we further executed \textbf{isolated clean-room Docker kernel builds} (\texttt{--network=none}, unprivileged \texttt{nobody}, read-only source mounts) plus programmatic \texttt{\#print axioms} (\texttt{Lean.collectAxioms}) across all \textbf{14} standalone compact repositories ($<15\text{k}$~LoC across T1 and T2): \textbf{11/14 (78.6\%)} build cleanly out of the box (auditing \textbf{3,190} theorems; median \textbf{28.4\,s}), with \textbf{0/14} genuine proof-compilation errors and \textbf{0/14} \texttt{Lean.ofReduceBool} dependencies. All three out-of-the-box failures are packaging/dependency omissions (two missing \texttt{lakefile}/\texttt{lean-toolchain} whose proofs pass cleanly in synthetic documented environments, bringing total audited theorems to \textbf{3,254}, plus one unbundled \texttt{../qanary-*} sibling path). Crucially, \textbf{7/8 (87.5\%)} compact \texttt{Tier3\_Gold} repositories build out of the box and \textbf{7/7 (100\%)} are transitively free of \texttt{sorryAx}, \texttt{Lean.ofReduceBool}, and custom axioms, whereas dynamic axiom auditing catches transitive \texttt{sorryAx} in \texttt{2611.16690} (9/570 theorems) and \textbf{11} custom axioms in \texttt{2605.16407}---both already flagged statically under Mode~\#6.
\end{abstract}

% -----------------------------------------------------------------------------
\section{Introduction and Motivation}
\label{sec:intro}

Interactive theorem proving in Lean~4~\cite{demoura2021lean4} has historically been concentrated in formal mathematics and programming-languages communities. Throughout 2026, however, autonomous LLM proving agents have lowered the barrier to generating large Lean~4 formalizations alongside research manuscripts. At the same time, high-profile audits of bulk AI-generated mathematical releases have demonstrated that a compiling Lean build is a weak proxy for trustworthy scholarship: semantic audits document silent hypothesis strengthening and intermediate lemma drift~\cite{bastounis2026navier,leanatlas2026}, multi-agent studies uncover emergent notation hijacking and verification-harness exploits~\cite{paglieri2026swarms}, and multi-million-line machine-generated codebases expose severe human-intelligibility and checker-scalability bottlenecks~\cite{axle2026}. In particular, Bastounis et al.~\cite{bastounis2026navier} distinguish \emph{Type~(i) autoformalization} (producing some compiling Lean proof inspired by a manuscript) from \emph{Type~(ii) autoformalization} (faithfully translating every definition, hypothesis, and lemma of the natural-language paper), showing that compiler-driven agents routinely mutate hypotheses or proof strategies to achieve a green build.

In this study, we ask five empirical questions about \emph{cutting-edge engineering, optimization, and systems research}:
\begin{enumerate}
  \item \textbf{Longitudinal Inflection:} How fast did Lean~4 adoption grow across engineering arXiv preprints between Nov~2025, May~2026, and Nov~2026 relative to traditional formal-methods tools?
  \item \textbf{Prevalence \& Disciplinary Distribution:} Where outside pure mathematics and traditional formal-methods venues are researchers actively shipping Lean proofs on arXiv right now, and how sharply does adoption vary across ECE and applied-theory subfields?
  \item \textbf{Authoring Paradigms:} Among papers contributing Lean proofs, which Lean$\leftrightarrow$\LaTeX{} workflow are authors actually following---post-hoc autoformalization, Lean-first authoring, synchronized blueprints, or spec-locked solving?
  \item \textbf{Verification Hygiene, Reproducibility \& Kernel Ground Truth:} What observable soundness, scope, transparency, and scale risks appear across public GitHub, GitLab, and Zenodo Lean artifacts, and do compact engineering repositories actually build cleanly in an isolated network-disabled container with zero transitive \texttt{sorryAx} or \texttt{Lean.ofReduceBool}?
  \item \textbf{Layer-A Semantic Fidelity:} When audited theorem-by-theorem against the paper's \LaTeX{} equations, do independent author-owned engineering Lean repos suffer the same Type~(i) statement drift documented in bulk AI autoformalization dumps, and how has fidelity evolved between May and Nov~2026?
\end{enumerate}

To ground domain selection in a concrete engineering lab profile, Pilot~v1 anchored its initial six categories to USC's Autonomous Networks Research Group (ANRG; 100 journal and 270 conference papers across IEEE/ACM venues such as ToN, TMC, JSAC, IoT-J, INFOCOM, MobiCom, SenSys, ISIT, WiOpt, IPDPS, and IROS): \texttt{cs.DC}, \texttt{cs.NI}, \texttt{cs.IT}, \texttt{eess.SY}, \texttt{cs.CR}, and \texttt{cs.RO}. Pilot~v2 expanded the exact same three-week window to \textbf{Signal Processing (\texttt{eess.SP})} and \textbf{Optimization and Control (\texttt{math.OC})}, added Zenodo and GitLab profiling plus Stage-2 negative recall calibration, and audited 17 Nov~2026 \texttt{engineering\_math\_proof} repos. Pilot~v0.4 added two matched historical baseline slices---\textbf{T0 (Nov~1--19, 2025; $n=2{,}435$)} and \textbf{T1 (May~1--19, 2026; $n=3{,}298$)}---bringing the total screened corpus to \textbf{9,837 unique preprints} and expanding our Phase~2 Layer-A audit cohort across all \textbf{23} locally inspected \texttt{engineering\_math\_proof} packages (20 in T2, including statically unpacked arXiv-ancillary and zipped archives, plus 3 in T1). Pilot~v0.5 adds \textbf{isolated Docker dynamic kernel verification} (\texttt{lake build} + programmatic \texttt{\#print axioms}) across all \textbf{14} compact standalone repositories ($<15\text{k}$~LoC across T1 and T2), elevating the study from purely static screening to a dynamic reproducibility benchmark.

% -----------------------------------------------------------------------------
\section{Background: Failure Modes, Solutions, and Authoring Paradigms}
\label{sec:taxonomy}

Because Lean's kernel only certifies that a formal proof term derives a formal proposition under a specific environment, a clean \texttt{lake build} leaves three distinct layers of scholarly trust exposed: \textbf{Layer~A} (\emph{semantic and statement-level fidelity}---``what was actually proved?''), \textbf{Layer~B} (\emph{adversarial and kernel-level integrity}---``was it legitimately proved?''), and \textbf{Layer~C} (\emph{structural, cognitive, and operational utility}---``is the artifact usable, auditable science?''). Synthesizing documented 2025--2026 formal-verification postmortems and tooling proposals~\cite{bastounis2026navier,zhu2026leanarchitect,paglieri2026swarms,leanatlas2026,axle2026,proofbridge2025}, Table~\ref{tab:taxonomy} organizes twelve failure modes across these three layers alongside their community defense-in-depth remedies and specifies which modes are observable in our study.

\begin{table}[!tbp]
\centering
\caption{Three-layer taxonomy of twelve failure modes in AI-assisted Lean~4 formalizations, corresponding community defense-in-depth remedies, and observability in this study (\textbf{P1} = Phase~1 static artifact profiling; \textbf{P2} = Phase~2 adversarial \LaTeX{}$\leftrightarrow$Lean audit; \textbf{Dyn.} = isolated Docker \texttt{lake build} + \texttt{\#print axioms} on compact repos $<15\text{k}$~LoC, Section~\ref{sec:sandbox}).}
\label{tab:taxonomy}
\scriptsize
\setlength{\tabcolsep}{2.8pt}
\renewcommand{\arraystretch}{1.04}
\begin{tabularx}{\textwidth}{@{}>{\raggedright\arraybackslash}p{2.35cm}>{\raggedright\arraybackslash}X>{\raggedright\arraybackslash}X>{\raggedright\arraybackslash}p{1.25cm}@{}}
\toprule
\textbf{Failure Mode} & \textbf{Mechanism \& Scholarly Risk} & \textbf{Defense-in-Depth Remedy} & \textbf{Scope} \\
\midrule
\multicolumn{4}{@{}l}{\textit{\textbf{Layer A: Semantic \& Statement-Level Fidelity (``What Was Actually Proved?'')}}} \\
\textbf{\#1 Statement / Proof Drift} & Type~(i) compilation silently strengthens hypotheses, weakens bounds, or alters intermediate lemmas relative to prose~\cite{bastounis2026navier}. & Bidirectional blueprints~\cite{zhu2026leanarchitect}; formal bidirectional equivalence checking~\cite{proofbridge2025}; cross-LLM semantic audits~\cite{bastounis2026navier} + human review. & \textbf{P2} \\
\textbf{\#2 Scope Downgrading} & Headline claims outstrip the Lean certificate (e.g., special cases, omitted complexity/asymptotics, or ``scope-note erasure''). & Explicit claim-level scope notes and verification manifests placed directly beside paper claims. & \textbf{P1 + P2} \\
\textbf{\#3 Definition Smuggling} & Custom definitions unanchored to Mathlib subtly distort standard mathematical objects, making target theorems artificially weak~\cite{leanatlas2026}. & Mathlib-first specification boundary; dependency-graph pruning of non-Mathlib defs (\texttt{Lean Atlas}~\cite{leanatlas2026}). & \textbf{P1 + P2} \\
\textbf{\#4 Junk Values \& Vacuity} & Totalized Lean operations (e.g., $1/0 = 0$, \texttt{Nat} subtraction $3-5=0$) or inconsistent/empty hypotheses yield vacuously true theorems. & Compilable non-vacuity witnesses (\texttt{example : $\exists x$, ...}), edge-case tests, and \texttt{EReal}/\texttt{ENNReal} anti-junk lifting. & \textbf{P2} \\
\midrule
\multicolumn{4}{@{}l}{\textit{\textbf{Layer B: Adversarial \& Kernel-Level Integrity (``Was It Legitimately Proved?'')}}} \\
\textbf{\#5 Reward Hacking / Notation} & Autonomous provers override notation, shadow identifiers, or tamper with targets to pass checks with zero \texttt{sorry}~\cite{paglieri2026swarms}. & Freeze trusted spec (\texttt{Challenge.lean}) vs.\ untrusted \texttt{Solution.lean} via sandboxed \texttt{lake comparator}; independent auditor agents. & \textbf{P1} (\texttt{P4} proxy) / Dyn. \\
\textbf{\#6 \texttt{sorry}, Custom \texttt{axiom}, \texttt{native\_decide}} & Transitive \texttt{sorryAx}, unproved \texttt{axiom}s, or \texttt{native\_decide} (\texttt{Lean.ofReduceBool}) bypass kernel reduction and expand the TCB to the compiler. & Programmatic \texttt{\#print axioms} restricted to \{\texttt{propext}, \texttt{Classical.choice}, \texttt{Quot.sound}\}; CI bans on \texttt{native\_decide}. & \textbf{P1 + Dyn.} \\
\textbf{\#7 Metaprog.\ \& Side Effects} & Elaborator/\texttt{.olean} tampering can bypass kernel checks; compile-time \texttt{\#eval} executes arbitrary host OS actions. & Network-isolated read-only sandboxes; static bans on \texttt{unsafe}, \texttt{implemented\_by}, and solution metaprogramming. & \textbf{P1 + Dyn.} \\
\textbf{\#8 Kernel Soundness Bugs} & Rare implementation bugs in Lean's C++ kernel (especially via metaprogramming) create a single-checker monoculture risk. & Multi-kernel replay (\texttt{lake check --paranoid}) across independent/verified checkers (e.g., \texttt{nanoda}, \texttt{lean4lean}). & Out of scope (Dyn.) \\
\textbf{\#9 Foundational Drift} & Lean's dependent type theory with an infinite universe hierarchy is proof-theoretically stronger than first-order ZFC. & Audit universe polymorphism and disclose extra foundational assumptions for logic/set-theory claims. & Out of scope (Dyn.) \\
\midrule
\multicolumn{4}{@{}l}{\textit{\textbf{Layer C: Structural, Cognitive, and Operational Utility (``Is It Usable, Auditable Science?'')}}} \\
\textbf{\#10 Proof Bloat / Tech.\ Debt} & Multi-10k to multi-million-LoC AI codebases produce truth certificates faster than recoverable conceptual structure. & Mandatory conceptual distillation/refactoring pass; separate statement interfaces from proof modules. & \textbf{P1} ($>20\text{k}$ LoC) \\
\textbf{\#11 Checker Timeouts} & External verification (\texttt{comparator}) hits high latency and timeouts at repo scale ($\approx 96\text{s}$ median, 6.2\% timeouts~\cite{axle2026}). & Tiered verification: fast inner-loop checkers (\texttt{Axle}~\cite{axle2026}) during search; incremental \texttt{comparator} at release. & \textbf{Dyn.} (build time) \\
\textbf{\#12 Process Opacity} & Missing/broken repos, unpinned toolchains, silent revisions, or mixed formalized/unformalized claims without clear provenance. & Standardized verification manifests (tier badge per claim, toolchain/Mathlib SHA, immutable releases). & \textbf{P1 + Dyn.} \\
\bottomrule
\end{tabularx}
\end{table}

Two technical distinctions in Table~\ref{tab:taxonomy} are especially relevant for engineering readers unfamiliar with Lean~4 internals. First, \textbf{Mode~\#6 groups \texttt{native\_decide} with \texttt{sorry} and custom \texttt{axiom}s} because \texttt{native\_decide} delegates finite checks to compiled native code rather than kernel reduction: this expands the trusted computing base (TCB) to the entire Lean compiler/runtime and prevents external checkers from verifying the proof, whereas pure kernel-checked classical Lean proofs depend exclusively on Lean's three standard axioms (\texttt{propext}, \texttt{Classical.choice}, \texttt{Quot.sound}). Second, \textbf{pinned specification checking via \texttt{lake comparator}} defeats Mode~\#5 and Mode~\#7 tampering by exporting proof terms from an untrusted \texttt{Solution.lean} inside a sandbox and matching them against a frozen \texttt{Challenge.lean}---yet it remains ``garbage-in, garbage-out'' if the specification itself omits hypotheses or conclusions, making independent Layer-A spec auditing indispensable (as we directly observe in Section~\ref{sec:phase2}).

\subsection{Lean\texorpdfstring{$\leftrightarrow$\LaTeX{}}{<->LaTeX} Authoring Paradigms and Trade-Offs}
How authors organize the interplay between natural-language exposition and Lean code strongly shapes which Layer-A and Layer-C failures arise. We classify every confirmed Lean-contributing paper into one of five authoring paradigms, whose structural trade-offs are summarized in Table~\ref{tab:paradigms}:
\begin{description}
  \item[\texttt{P1} (Paper-First, Post-Hoc Autoformalization):] Informal proofs and the \LaTeX{} manuscript are written first; Lean code is created afterward (often by AI agents) and linked via footnotes, manual tables, SHA-256 hash locks, or post-hoc correspondence maps without bidirectional sync tooling. \emph{Trade-off:} preserves natural mathematical discovery, but risks silent Type~(i) drift whenever a formalizer patches Lean hypotheses without updating the manuscript.
  \item[\texttt{P2} (Raw Lean-First / Auto-Informalization):] The kernel-checked Lean proof is built first, and the human manuscript is authored or decompiled directly against the Lean dependency graph. \emph{Trade-off:} prevents downward lemma drift by construction, but does not guard against an erroneous initial top-level specification and risks propagating machine-generated proof bloat and type-theoretic plumbing into prose.
  \item[\texttt{P3} (Spec-Pinned Synchronized Co-Development):] Bidirectional literate tooling such as \texttt{LeanArchitect} (\texttt{@[blueprint]}, \texttt{\textbackslash inputleannode})~\cite{zhu2026leanarchitect}, \texttt{leanblueprint}, or \texttt{Verso Blueprints} keeps the \LaTeX{} theorem DAG and Lean declarations structurally synchronized, automatically inferring dependencies and \texttt{sorry}-free status while encapsulating low-level tactic plumbing inside each node.
  \item[\texttt{P4} (Spec-Lock $\rightarrow$ AI-Solve $\rightarrow$ Distill):] Target statements or formal optimization setups are frozen in a read-only specification module (e.g., \texttt{Challenge.lean}, \texttt{ReferenceClaims}, or \texttt{LongTails/Setups}) checked against untrusted proofs via \texttt{lake comparator} and external kernels (e.g., \texttt{nanoda}). Separating the spec owner from the solver prevents self-grading and reward hacking during unconstrained AI search, followed by conceptual distillation into clean exposition.
  \item[\texttt{P0} (Opaque / Unspecified):] The paper claims Lean verification of its results but supplies neither inspectable Lean artifacts nor sufficient workflow description to classify the paradigm.
\end{description}

\begin{table}[!tbp]
\centering
\caption{Head-to-head trade-offs of the four inspectable Lean$\leftrightarrow$\LaTeX{} authoring paradigms (\texttt{P0} denotes opaque/unspecified workflows).}
\label{tab:paradigms}
\scriptsize
\setlength{\tabcolsep}{3.5pt}
\begin{tabularx}{\textwidth}{@{}>{\raggedright\arraybackslash}p{3.0cm}>{\raggedright\arraybackslash}X>{\raggedright\arraybackslash}X>{\raggedright\arraybackslash}X>{\raggedright\arraybackslash}X@{}}
\toprule
\textbf{Paradigm} & \textbf{Top-Level Spec Fidelity} & \textbf{Intermediate Lemma Sync} & \textbf{Readability / Low Debt} & \textbf{Tooling \& Workflow Overhead} \\
\midrule
\texttt{P1} Paper-First (Post-Hoc) & Low--Moderate (drift risk) & Poor (Type~(i) drift prone) & High prose quality (unverified) & Low upfront; high post-hoc audit cost \\
\texttt{P2} Raw Lean-First & Moderate (bad initial spec risk) & Very High (1-to-1 by construction) & Poor--Moderate (plumbing/bloat) & Moderate (early formal rigidity) \\
\texttt{P3} Synchronized Blueprint & High (pinned theorem DAG) & Very High (auto-inferred DAG) & Very High (tactics encapsulated) & Moderate--High (blueprint tooling) \\
\texttt{P4} Spec-Lock + Solve + Distill & Very High (frozen \texttt{Challenge}) & Very High (if distilled to prose) & High (with distillation pass) & High (spec audit + \texttt{comparator} CI) \\
\bottomrule
\end{tabularx}
\end{table}

\subsection{Observed Static Verification Tiers, Domain Roles, and Risk Flags}
In a full publication-grade verification pipeline, gold-standard certification requires human-audited specifications, sandboxed \texttt{comparator} plus multi-kernel replay, and synchronized blueprints. Because our study strictly avoids executing untrusted build scripts or compile-time metaprograms (Mode~\#7), we operationalize verification maturity via \emph{statically observed} tiers---treating them as an observable proxy that Phase~2 then stress-tests semantically on Layer~A:
\begin{itemize}
  \item \textbf{Domain Role:} \texttt{engineering\_math\_proof} (Lean verifies mathematical theorems in optimization, control, system ID, coding/information theory, signal processing, statistical learning, or cryptographic math), \texttt{software\_or\_system\_fv} (Lean verifies an implementation, policy engine, CDC pipeline, or security calculus), or \texttt{pure\_tcs\_or\_math} (theoretical CS or pure math cross-listed into engineering).
  \item \textbf{Observed Verification Tier (Static):} \texttt{Tier3\_Gold} (main results publicly formalized, zero \texttt{sorry}/custom \texttt{axiom}/\texttt{native\_decide} in proved endpoints, plus explicit author-supplied audit evidence such as \texttt{\#print axioms}, \texttt{comparator} scripts, hash-locks, or external-kernel replay); \texttt{Tier2\_Silver} (public compiling-looking Lean repo covering supporting lemmas/partial scope or lacking formal audit evidence); \texttt{Tier1\_Bronze} (no inspectable public repo, unresolved external axioms, or opaque process).
  \item \textbf{Static Risk Flags (Phase~1) \& Semantic Layer-A Modes (Phase~2):} Phase~1 tracks the six statically detectable modes from Table~\ref{tab:taxonomy}: \texttt{Mode2\_ScopeDowngrade}, \texttt{Mode3\_CustomDefs}, \texttt{Mode6\_SorryOrAxiomOrNativeDecide}, \texttt{Mode7\_MetaprogOrUnsafe}, \texttt{Mode10\_ProofBloat} ($>20\text{k}$ LoC), and \texttt{Mode12\_ProcessOpacity}. Phase~2 executes theorem-by-theorem \LaTeX{}$\leftrightarrow$Lean comparison across all four Layer-A modes (Modes~\#1--\#4; Section~\ref{sec:phase2}).
\end{itemize}

% -----------------------------------------------------------------------------
\section{Reproducible Tiered Measurement Pipeline \& Recall Calibration}
\label{sec:pipeline}

Because Lean adoption in non-formal-methods engineering is a rare event ($\approx 1\%$ base rate), running an LLM over every paper wastes token budget, whereas title/abstract filtering misses papers mentioning Lean only in later sections, footnotes, or acknowledgments. We built a five-stage, SQLite-checkpointed, idempotent pipeline:

\begin{enumerate}
  \item \textbf{Stage 1 --- arXiv Metadata Ingest (0 LLM tokens):} Query arXiv via Atom API (Pilot~v1) and arXiv's official OAI-PMH \texttt{arXivRaw} harvester (\texttt{oaipmh.arxiv.org/oai}, Pilot~v2) for v1 submissions in \texttt{2026-11-01}--\texttt{2026-11-22} across all 8 target categories, deduplicating cross-listed preprints with set-unioned target tags in \texttt{data/survey.db}. This yielded \textbf{4,104} unique preprints (latest announced submission: November~19, 2026; \textbf{4,060} excluding pure-TCS primary categories \texttt{cs.CC}, \texttt{cs.DS}, \texttt{cs.LO}, \texttt{cs.DM}, \texttt{cs.FL}, \texttt{cs.CG}). Metadata-only screening found only \textbf{16} unambiguous Lean hits.
  \item \textbf{Stage 2 --- Polite Full-Text Pre-Screen (0 LLM tokens):} Fetch \texttt{arxiv.org/html/\{id\}v\{v\}} first, falling back to \texttt{arxiv.org/pdf/\{id\}v\{v\}} via \texttt{pypdf} on 404. We achieved \textbf{99.83\% full-text coverage} (3,849 HTML5, 248 PDF, 7 genuine HTTP 404s). Deterministic case-aware regex detectors scanned full text for strong Lean indicators (\texttt{Lean 4}, \texttt{Mathlib}, \texttt{lakefile}, \texttt{LeanArchitect}, \texttt{Verso}, \texttt{sorryAx}, \texttt{native\_decide}, etc.) and contextual \texttt{Lean} mentions near proof terminology (excluding phrases like ``lean manufacturing''). Full-text scanning surfaced \textbf{55 Lean candidates} (47 strong, 8 contextual-only---$3.4\times$ metadata-only recall), eliminating \textbf{98.7\%} of the corpus before any LLM call.
  \item \textbf{Stage 3 --- Strictly Static Artifact Profiler (0 LLM tokens):} For all 55 Lean candidates, download raw \LaTeX{} source (\texttt{/src/\{id\}}) to detect blueprint macros, theorem counts, AI disclosures, embedded artifact URLs, and bundled arXiv ancillary Lean trees/archives. For linked public repositories across \textbf{GitHub} (authenticated \texttt{gh api} + shallow \texttt{git clone --depth 1} up to 600\,MB, including safe static unpacking of zipped Lean archives), \textbf{GitLab} (including self-hosted institutional instances), and \textbf{Zenodo} (\texttt{zenodo.org/api/records} + safe static archive extraction), run a comment-stripping Lean lexer without ever executing \texttt{lake build} or \texttt{\#eval}.
  \item \textbf{Stage 4 --- Parallel Lightweight-Agent Classification:} Using Google Antigravity's dynamic \texttt{run\_workflow} orchestrator with strict JSON schema validation, read-only \texttt{research} agents at the \texttt{flash} tier (\textbf{Gemini~4~Hydrogen, High effort}) classified all 55 Lean candidates (inspecting local READMEs, lakefiles, audit manifests, and representative \texttt{.lean} files). Pilot~v1 also classified 34 Beyond-Lean candidates at the \texttt{flash\_lite} tier (\textbf{Gemini~3.5~Flash-Lite}) across the original 6 categories; Pilot~v2 focused LLM budget exclusively on Lean.
  \item \textbf{Stage 5 (Phase~2) --- Adversarial Layer-A Semantic Fidelity Audit:} For all \textbf{20} confirmed T2 \texttt{engineering\_math\_proof} papers with locally inspected Lean repositories (including statically unpacked arXiv ancillary and zipped archives), we extracted math-preserving Markdown from arXiv HTML5 (\texttt{<math alttext="...">}, with PDF fallback) and fanned out parallel \texttt{pro}-tier (\textbf{Gemini~4~Hydrogen, Extra High effort}) read-only adversarial referee agents (building on the cross-LLM semantic-auditing observation of~\cite{bastounis2026navier}) pairing 2--5 principal paper theorems against their exact Lean declarations and transitive custom definitions with \texttt{file:line} citations.
\end{enumerate}

\subsection{Stage-2 Recall Calibration on Negatives}
\label{sec:recall}
A central methodological threat in regex-gated empirical studies is false negatives---specifically, papers that link a Lean repository without naming Lean in their prose. Using \texttt{src/calibrate\_recall.py} (seed \texttt{20261122}), we partitioned all \textbf{4,042} full-text-screened Stage-2 negatives into 1,390 papers containing non-boilerplate GitHub links (34.4\%) and 2,652 without. We audited a stratified random sample of \textbf{200 GitHub-linked Stage-2 negatives} (querying recursive file trees for \textbf{333 unique GitHub repositories} via \texttt{gh api}; 195 papers had $\ge 1$ accessible public repo) for any \texttt{*.lean}, \texttt{lean-toolchain}, or \texttt{lakefile.*} file, plus \textbf{100 random non-GitHub negatives} for broad metadata/comment proof keywords.

\textbf{Result:} We detected \textbf{0 / 200} missed Lean repositories among GitHub-linked negatives and \textbf{0 / 100} broad metadata flags among non-GitHub negatives (archived in \texttt{reports/slices/2026-11/recall\_calibration.md}). This yields an exact one-sided 95\% Clopper--Pearson upper bound of \textbf{1.49\%} on hidden-Lean prevalence among GitHub-linked negatives (\textbf{1.52\%} conditioning on accessible repos), confirming high Stage-2 recall. Figure~\ref{fig:pipeline} summarizes the complete five-stage funnel and recall-calibration design.

\begin{figure}[t]
\centering
\includegraphics[width=\textwidth]{fig_pipeline_funnel.pdf}
\caption{Five-stage, SQLite-checkpointed measurement pipeline and Stage-2 negative recall calibration. Stages~1--3 and recall calibration consume zero LLM tokens, pruning 98.7\% of the 4,104-paper corpus before parallel \texttt{flash} classification and \texttt{pro} Layer-A auditing.}
\label{fig:pipeline}
\end{figure}

% -----------------------------------------------------------------------------
\section{Empirical Results: Prevalence, Paradigms, and Static Hygiene}
\label{sec:results}

\subsection{Disciplinary Bifurcation Across Eight Engineering Categories}
Across all 4,104 preprints, Stage~4 confirmed \textbf{39 papers contributing Lean proofs} (0.95\% of the full corpus; 13 of the 55 candidates were citations/discussions and 3 were false positives). Removing papers whose primary category is pure theoretical CS (\texttt{cs.CC}, \texttt{cs.DS}, \texttt{cs.LO}, \texttt{cs.DM}, \texttt{cs.FL}, \texttt{cs.CG}) leaves \textbf{36 confirmed Lean papers out of 4,060 (0.89\%)}. Crucially, \textbf{three} of the 8 contextual-only Lean hits (\texttt{2611.15431}, \texttt{2611.20131}, and \texttt{2611.02169}) proved to be genuine Lean papers, validating contextual full-text matching.

Table~\ref{tab:categories} and Figure~\ref{fig:categories} reveal a stark disciplinary split across engineering fields:
\begin{itemize}
  \item \textbf{Optimization \& Control (\texttt{math.OC}) is the epicenter of engineering Lean adoption:} \textbf{19 / 724 papers (2.62\%)} contribute Lean proofs, of which \textbf{17 (89.5\%)} are \texttt{engineering\_math\_proof} theorems (heavy-ball momentum, ADMM, alternating mirror descent, Wolfe's min-norm-point method, McKean--Vlasov SDE regression, discrete modified-gradient systems, and oracle-complexity autoresearch).
  \item \textbf{Information Theory (\texttt{cs.IT}, 1.78\%), Distributed Computing (\texttt{cs.DC}, 1.79\%), Crypto/Security (\texttt{cs.CR}, 1.11\%), and Systems/Control (\texttt{eess.SY}, 0.41\%)} form a strong secondary tier. In \texttt{cs.IT}, 6/6 Lean papers are \texttt{engineering\_math\_proof} (vs.\ 0 Beyond-Lean papers), whereas \texttt{cs.DC} Lean papers are predominantly theoretical distributed algorithms (\texttt{pure\_tcs\_or\_math}, 5/6).
  \item \textbf{Zero adoption in three major applied communities (0 / 1,817 papers):} Signal Processing (\texttt{eess.SP}, 0/462), Networking (\texttt{cs.NI}, 0/196), and Robotics (\texttt{cs.RO}, 0/1,159) recorded \textbf{zero} confirmed Lean papers across the entire three-week window.
\end{itemize}

\begin{table}[t]
\centering
\caption{Three-week Lean prevalence across eight engineering and applied-theory arXiv categories (Nov~1--19, 2026). Category counts include cross-listings; the bottom rows report the deduplicated union and the non-TCS-primary engineering subset. Beyond-Lean LLM baseline was run on the six Pilot~v1 categories.}
\label{tab:categories}
\footnotesize
\setlength{\tabcolsep}{3.5pt}
\begin{tabular}{lrrrrr}
\toprule
\textbf{Category} & \textbf{Papers} & \textbf{Lean Contrib.} & \textbf{Lean (\%)} & \textbf{Eng.\ Math Proof} & \textbf{Beyond-Lean (v1)} \\
\midrule
\texttt{math.OC} (Optimization \& Control) & 724 & \textbf{19} & \textbf{2.62\%} & \textbf{17} & \textemdash{} \\
\texttt{cs.DC} (Distributed Computing) & 335 & 6 & 1.79\% & 0 & 3 \\
\texttt{cs.IT} (Information Theory) & 338 & 6 & 1.78\% & 6 & 0 \\
\texttt{cs.CR} (Crypto \& Security) & 629 & 7 & 1.11\% & 4 & 8 \\
\texttt{eess.SY} (Systems \& Control) & 738 & 3 & 0.41\% & 2 & 1 \\
\texttt{eess.SP} (Signal Processing) & 462 & 0 & 0.00\% & 0 & \textemdash{} \\
\texttt{cs.NI} (Networking) & 196 & 0 & 0.00\% & 0 & 0 \\
\texttt{cs.RO} (Robotics) & 1,159 & 0 & 0.00\% & 0 & 2 \\
\midrule
\textbf{Deduplicated Union (8 cats)} & \textbf{4,104} & \textbf{39} & \textbf{0.95\%} & \textbf{28} & \textbf{12} \\
\emph{Excl.\ pure-TCS primary cats} & \emph{4,060} & \emph{36} & \emph{0.89\%} & \emph{27} & \emph{\textemdash{}} \\
\bottomrule
\end{tabular}
\end{table}

\begin{figure}[t]
\centering
\includegraphics[width=0.92\textwidth]{fig_category_bifurcation.pdf}
\caption{Confirmed Lean~4 adoption rate (\%) across all eight engineering arXiv categories (Nov~1--19, 2026), stacked by domain role. Optimization and Control (\texttt{math.OC}, 2.62\%) leads the study, whereas Signal Processing (\texttt{eess.SP}), Networking (\texttt{cs.NI}), and Robotics (\texttt{cs.RO}) show zero adoption across 1,817 papers.}
\label{fig:categories}
\end{figure}

\subsection{Authoring Paradigms, AI Assistance, and Artifact Completeness}
Table~\ref{tab:distributions} summarizes classifications across all 39 confirmed Lean papers. Five empirical patterns stand out:

\begin{enumerate}
  \item \textbf{Post-hoc \texttt{P1} dominates (71.8\%); synchronized blueprints (\texttt{P3}) and Lean-first (\texttt{P2}) remain at 0\%.} Even after nearly doubling our sample from 20 to 39 papers, \emph{not a single engineering paper} used \texttt{LeanArchitect}, \texttt{leanblueprint}, or \texttt{Verso Blueprints}. Instead, 28 papers formalized prose theorems post-hoc, with several independently inventing lightweight sync guards---most notably \texttt{2611.18455} (\texttt{inody/wmfki\_lean}), which hashes manuscript \LaTeX{} theorem environments in \texttt{statement-lock.json}, and \texttt{2611.15201} (Zenodo), which wraps a post-hoc PDF formalization inside a \texttt{Challenge.lean}/\texttt{Solution.lean} \texttt{comparator} harness.
  \item \textbf{Three flagship \texttt{P4} spec-lock adopters (7.7\%):} Alongside \texttt{2611.14475} (explicit RIP matrices) and \texttt{2611.04876} (matrix refreshing), \texttt{math.OC} adds \textbf{\texttt{2611.13872} (\emph{LongTails})}, where humans freeze formal optimization setups in Lean~4, AI agents discover algorithms and prove oracle complexity bounds, and CI (\texttt{verify.yml}) dynamically generates \texttt{Challenge} modules checked by \texttt{leanprover/comparator}.
  \item \textbf{AI-assisted formalization is the norm (76.9\% explicit, 0\% explicitly human-only):} Thirty papers explicitly credit LLM agents (Claude Opus~5.5/Fable~5.1/Claude Code, GPT-6 Astra/GPT-6.1 Sol/Codex, DeepSeek, Gemini, Harmonic); the remaining 9 leave authorship method unstated.
  \item \textbf{Zenodo, large-monorepo, and ancillary/zip archive profiling closes Stage-3 blind spots:} Extending Stage~3 to Zenodo records and 600\,MB GitHub repos retrieved \texttt{2611.15201} (Zenodo~24127514: 5.3k~LoC, sorry-free \texttt{Solution.lean} + \texttt{comparator}), \texttt{2611.10810} (Zenodo~24035424: 10.9k~LoC, 475 \texttt{\#print axioms}), \texttt{2611.15128} (\texttt{leanEthereum/leanVM}: 20 \texttt{AffineCombiner*.lean} files, 5.0k~LoC, CI + \texttt{\#print axioms}), and \texttt{2611.12680} (Zenodo~24136107: 23.9k~LoC, zero sorry), upgrading all three previously uninspected v1 papers to \texttt{Tier3\_Gold}. Unpacking bundled arXiv ancillary packages and GitHub zip archives further recovered \texttt{2611.26140} (28.0k~LoC, 53 files, 0 \texttt{sorry}/axioms), \texttt{2611.16781} (34.6k~LoC, 121 files, 1,192 \texttt{\#print axioms}, 0 \texttt{sorry}/axioms), and \texttt{2611.12585} (60.7k~LoC, 341 files, 2,715 \texttt{\#print axioms}, 0 \texttt{sorry}/axioms), leaving only \texttt{2611.22732}'s institutional GitLab host (\texttt{gitlab.mpi-sws.org}, persistent port-443 timeout on retry) and \texttt{2611.09169} (external portal) unprofiled alongside 9 papers lacking a public repository (6 none, 3 broken/empty links); frozen initial Stage-4 classifications keep \texttt{Mode12\_ProcessOpacity} at 43.6\% (17/39).
  \item \textbf{Over four million lines of Lean in three weeks:} Across 27 unique locally scanned primary artifacts (covering 28 papers; Figure~\ref{fig:loc}), Lean~4 codebases total \textbf{4.30~million non-empty lines of code}, spanning compact 528--2,191~LoC developments up to multi-paper programs exceeding 697k--929k~LoC.
\end{enumerate}

\begin{table}[t]
\centering
\caption{Classification distributions across all 39 confirmed Lean-contributing papers (Nov~1--19, 2026).}
\label{tab:distributions}
\small
\begin{tabularx}{\textwidth}{lXr}
\toprule
\textbf{Dimension} & \textbf{Category / Label} & \textbf{Count (\%)} \\
\midrule
\textbf{Authoring Paradigm} & \texttt{P1\_paper\_first\_posthoc} (post-hoc formalization) & 28 (71.8\%) \\
 & \texttt{P0\_opaque\_or\_unspecified} (no inspectable repo / opaque workflow) & 8 (20.5\%) \\
 & \texttt{P4\_speclock\_solve\_distill} (\texttt{comparator} spec-lock) & 3 (7.7\%) \\
 & \texttt{P2\_lean\_first} / \texttt{P3\_blueprint\_sync} & 0 (0.0\%) \\
\midrule
\textbf{Domain Role} & \texttt{engineering\_math\_proof} (opt., control, IT, coding, stat.\ ML, crypto) & 28 (71.8\%) \\
 & \texttt{pure\_tcs\_or\_math} (CONGEST, LTL synthesis, H17, graphic TSP) & 7 (17.9\%) \\
 & \texttt{software\_or\_system\_fv} (policy engines, CDC, type-confusion, payment) & 4 (10.3\%) \\
\midrule
\textbf{Verification Tier} & \texttt{Tier3\_Gold} (public repo, clean endpoints, explicit axiom/kernel audit) & 19 (48.7\%) \\
(Static Observed) & \texttt{Tier2\_Silver} (public repo, partial scope or basic build) & 7 (17.9\%) \\
 & \texttt{Tier1\_Bronze} (no inspectable repo, custom axioms, or opaque process) & 13 (33.3\%) \\
\midrule
\textbf{AI Assistance} & \texttt{explicitly\_yes} / \texttt{unclear} / \texttt{explicitly\_no} & 30 / 9 / 0 \\
\textbf{Artifact Status} & Public repo (29: 27 scanned, 2 ext) / Ancillary (1, scanned) / Broken (3) / None (6) & 29 / 1 / 3 / 6 \\
\midrule
\textbf{Static Risk Flags} & \texttt{Mode12\_ProcessOpacity} (no repo, missing CI/toolchain, or unclear scope) & 17 (43.6\%) \\
(multi-label) & \texttt{Mode10\_ProofBloat} ($>20\text{k}$ LoC AI-scale codebase) & 12 (30.8\%) \\
 & \texttt{Mode2\_ScopeDowngrade} (partial lemmas, $K=2$, or omitted complexity) & 10 (25.6\%) \\
 & \texttt{Mode6\_SorryOrAxiomOrNativeDecide} (unresolved axioms / \texttt{sorry} / \texttt{native\_decide}) & 3 (7.7\%) \\
 & \texttt{Mode3\_CustomDefs} (heavy custom definitions unanchored to Mathlib) & 1 (2.6\%) \\
\bottomrule
\end{tabularx}
\end{table}

\begin{figure}[t]
\centering
\includegraphics[width=\textwidth]{fig_paradigms_and_loc.pdf}
\caption{(a) Distribution of Lean$\leftrightarrow$\LaTeX{} authoring paradigms stacked by observed verification tier across all 39 confirmed Lean papers (\texttt{P2} and \texttt{P3} remain at 0\%). (b) Non-empty Lean~4 lines of code (log scale) across all 28 confirmed papers with locally scanned GitHub, Zenodo, or arXiv-ancillary artifacts (*paper-specific subpackage LoC for multi-paper monorepos \texttt{PolynomialRealizationsProgram} and \texttt{leanVM}).}
\label{fig:loc}
\end{figure}

\subsection{Complete Inventory of Confirmed Lean Papers (39)}
Table~\ref{tab:inventory} catalogs all 39 confirmed Lean-contributing papers. Beyond the control/system-ID (\texttt{eess.SY}) and coding/information-theory (\texttt{cs.IT}) clusters identified in Pilot~v1, the Pilot~v2 \texttt{math.OC} cohort highlights a vibrant optimization ecosystem spanning first-order momentum (\texttt{2611.13584}, \texttt{2611.09050}), game-theoretic online learning (\texttt{2611.05243}), simplex/Wolfe geometry (\texttt{2611.12680}), stochastic particle SDE identification (\texttt{2611.18455}), dynamical-systems stability (\texttt{2611.14756}), automated oracle-complexity discovery (\texttt{2611.13872}, \texttt{2611.23080}), and bulk optimization autoformalization (\texttt{2611.26272}).

\begin{table}[p]
\centering
\caption{Complete inventory of all 39 confirmed Lean-contributing papers across the 4,104-paper T2 corpus (Nov~1--19, 2026). Role: \textbf{Eng} = \texttt{engineering\_math\_proof}, \textbf{Sys} = \texttt{software\_or\_system\_fv}, \textbf{TCS} = \texttt{pure\_tcs\_or\_math}. LoC reports non-empty lines in the paper's primary scanned artifact (\textemdash{} = none/broken; \emph{ext} = unreachable GitLab/external portal; * = paper-specific subpackage in multi-paper repo). Risks: \textbf{M2} Scope, \textbf{M3} CustomDefs, \textbf{M6} Axiom/\texttt{sorry}/\texttt{native\_decide}, \textbf{M10} Bloat, \textbf{M12} Opacity.}
\label{tab:inventory}
\tiny
\setlength{\tabcolsep}{2.3pt}
\begin{tabularx}{\textwidth}{llllllrlX}
\toprule
\textbf{arXiv ID} & \textbf{Prim.} & \textbf{Role} & \textbf{Par.} & \textbf{Tier} & \textbf{AI} & \textbf{LoC} & \textbf{Risks} & \textbf{Short Title \& Citation} \\
\midrule
2611.26272 & \texttt{math.OC} & Eng & \texttt{P1} & T2 & Yes & 151,214 & M10,M12 & Paper2Lean: Autoformalizing Math Papers (\texttt{OptPaper-20})~\cite{arxiv2611_26272} \\
2611.26140 & \texttt{math.OC} & Eng & \texttt{P1} & T1 & Yes & 28,043 & M2,M12 & Optimal Patrol Widths in Corridor Barrier-Guarding~\cite{arxiv2611_26140} \\
2611.25599 & \texttt{cs.CR} & Sys & \texttt{P0} & T1 & Unc & \textemdash{} & M12 & Tether: Verified Stateful Policy Engine~\cite{arxiv2611_25599} \\
2611.23762 & \texttt{quant-ph}* & Eng & \texttt{P1} & T2 & Yes & 589 & M2,M12 & Transversal CCZ Gates on qLDPC Codes~\cite{arxiv2611_23762} \\
2611.23080 & \texttt{cs.AI}* & Eng & \texttt{P0} & T1 & Yes & \textemdash{} & M12 & Colombo: Autoresearch for Oracle-Based Optimization~\cite{arxiv2611_23080} \\
2611.22732 & \texttt{cs.CR} & Eng & \texttt{P1} & T1 & Yes & \emph{ext} & M12 & Mixing Properties of Parallel Mixnets (GitLab timeout)~\cite{arxiv2611_22732} \\
2611.22468 & \texttt{math.OC} & Eng & \texttt{P0} & T1 & Yes & \textemdash{} & M12 & 5/2 Guarantee for Echelon Stock $(s,S)$ Policies~\cite{arxiv2611_22468} \\
2611.20298 & \texttt{cs.PL}* & Sys & \texttt{P1} & T1 & Yes & \textemdash{} & M12 & Type-Confusion Security~\cite{arxiv2611_20298} \\
2611.20131 & \texttt{math.NA}* & Eng & \texttt{P0} & T1 & Yes & \textemdash{} & M12 & Quadrature-Induced Local Minima (Lotka--Volterra)~\cite{arxiv2611_20131} \\
2611.19551 & \texttt{cs.CR} & Sys & \texttt{P1} & T3 & Yes & 1,332 & none & TrustBound: Payment Epochs for AI Agents~\cite{arxiv2611_19551} \\
2611.18455 & \texttt{stat.ML}* & Eng & \texttt{P1} & T3 & Unc & 12,236 & none & Mesh-Free Weak-Form Regression for Particle SDEs~\cite{arxiv2611_18455} \\
2611.18220 & \texttt{eess.SY} & Eng & \texttt{P1} & T3 & Unc & 528 & M12 & Lean for System Identification~\cite{arxiv2611_18220} \\
2611.16781 & \texttt{math.OC} & Eng & \texttt{P1} & T2 & Unc & 34,579 & M12 & Boundary Actuation of a Chain Fountain~\cite{arxiv2611_16781} \\
2611.16690 & \texttt{math.AG}* & TCS & \texttt{P1} & T2 & Yes & 6,592 & M2,M6 & Hilbert's 17th Problem for Symmetric Functions~\cite{arxiv2611_16690} \\
2611.15431 & \texttt{quant-ph}* & Eng & \texttt{P0} & T1 & Yes & \textemdash{} & M12 & Subadditivity in Nichols--Wootters Entropies~\cite{arxiv2611_15431} \\
2611.15201 & \texttt{cs.DC} & TCS & \texttt{P1} & T3 & Yes & 5,307 & M2 & Sparse Matrix Mult.\ Exp.\ 1.5 (Zenodo \texttt{comparator})~\cite{arxiv2611_15201} \\
2611.15128 & \texttt{cs.IT} & Eng & \texttt{P1} & T3 & Yes & 4,963* & M2 & Affine Combiners for RS Proximity Gaps (\texttt{leanVM})~\cite{arxiv2611_15128} \\
2611.14756 & \texttt{math.DS}* & Eng & \texttt{P1} & T3 & Unc & 670 & none & Stability of Discrete-Time Modified-Gradient Systems~\cite{arxiv2611_14756} \\
2611.14475 & \texttt{math.FA}* & Eng & \texttt{P4} & T3 & Yes & 29,967 & M2,M10 & Explicit Restricted Isometries $k=m^{2/3-o(1)}$~\cite{arxiv2611_14475} \\
2611.13872 & \texttt{math.OC} & Eng & \texttt{P4} & T3 & Yes & 697,377 & none & LongTails: Automated Research in Optimization~\cite{arxiv2611_13872} \\
2611.13585 & \texttt{cs.DS}* & TCS & \texttt{P1} & T1 & Unc & \textemdash{} & M12 & Approximating Graphic TSP by $b$-Matchings~\cite{arxiv2611_13585} \\
2611.13584 & \texttt{math.OC} & Eng & \texttt{P1} & T3 & Unc & 55,300* & M10 & Exact (Non-)Accelerations of Heavy-Ball Discretizations~\cite{arxiv2611_13584} \\
2611.12680 & \texttt{math.OC} & Eng & \texttt{P1} & T3 & Yes & 23,932 & M10 & Wolfe's Min-Norm-Point Generalizes Simplex (Zenodo)~\cite{arxiv2611_12680} \\
2611.12585 & \texttt{math.NA}* & Eng & \texttt{P1} & T2 & Yes & 60,664 & none & Gap in Galerkin Errors for KdV Equation (ancillary)~\cite{arxiv2611_12585} \\
2611.11713 & \texttt{cs.IT} & Eng & \texttt{P1} & T3 & Yes & 69,875 & M10 & Binary Belief Propagation Fixed Points~\cite{arxiv2611_11713} \\
2611.10902 & \texttt{eess.SY} & Eng & \texttt{P1} & T3 & Unc & 2,068 & M2 & Bisect-Index-Choose EV Charging~\cite{arxiv2611_10902} \\
2611.10810 & \texttt{cs.DB}* & Sys & \texttt{P1} & T3 & Yes & 10,862 & M3 & DBLog Incremental Snapshots (Zenodo, 475 ax.\ checks)~\cite{arxiv2611_10810} \\
2611.09169 & \texttt{math.OC} & Eng & \texttt{P0} & T1 & Yes & \emph{ext} & M12 & Multiplier Step Size in Two-Block ADMM (ReasBook)~\cite{arxiv2611_09169} \\
2611.09050 & \texttt{math.OC} & Eng & \texttt{P1} & T3 & Yes & 139,600* & M10 & Resolvent-Diagonal Energy of Isospectral Jacobi Matrices~\cite{arxiv2611_09050} \\
2611.08172 & \texttt{cs.IT} & Eng & \texttt{P1} & T3 & Yes & 200,335 & M10 & Three-Fold Product Expansion of AG Codes~\cite{arxiv2611_08172} \\
2611.05243 & \texttt{cs.LG}* & Eng & \texttt{P1} & T3 & Yes & 11,204 & none & Logarithmic Regret for Alternating Mirror Descent~\cite{arxiv2611_05243} \\
2611.05136 & \texttt{cs.CC}* & Eng & \texttt{P1} & T2 & Yes & 219,688 & M2,M10 & $TC^0$ Pseudorandomness from Lattices~\cite{arxiv2611_05136} \\
2611.04876 & \texttt{cs.CR} & Eng & \texttt{P4} & T3 & Yes & 81,481 & M10 & Leakage-Resilient Comp.\ via Matrix Refreshing~\cite{arxiv2611_04876} \\
2611.04561 & \texttt{cs.DC} & TCS & \texttt{P1} & T3 & Yes & 929,448 & M6,M10 & Distributed LTL Synthesis (Msg Adversaries)~\cite{arxiv2611_04561} \\
2611.02256 & \texttt{cs.DC} & TCS & \texttt{P1} & T2 & Yes & 2,191 & M2 & Planar Shortest Paths \& Max-Flow (CONGEST)~\cite{arxiv2611_02256} \\
2611.02169 & \texttt{math.OC} & Eng & \texttt{P0} & T1 & Yes & \textemdash{} & M12 & Split Cuts in Fixed Integer Dimensions~\cite{arxiv2611_02169} \\
2611.02057 & \texttt{cs.DS}* & TCS & \texttt{P1} & T1 & Yes & 166,354 & M6,M10,M12 & Covering Problems with Short Broadcasts~\cite{arxiv2611_02057} \\
2611.01761 & \texttt{cs.DC} & TCS & \texttt{P1} & T3 & Yes & 158,351 & M10 & Anonymous Directed Dynamic Networks~\cite{arxiv2611_01761} \\
2611.01384 & \texttt{math.OC} & Eng & \texttt{P0} & T1 & Unc & \textemdash{} & M2,M12 & Benign Population Landscapes Do Not Transfer to ERM~\cite{arxiv2611_01384} \\
\bottomrule
\multicolumn{9}{l}{\footnotesize *Cross-listed into one of the eight target categories (\texttt{math.OC}, \texttt{cs.IT}, \texttt{cs.CR}, \texttt{cs.DC}, or \texttt{eess.SY}).}
\end{tabularx}
\end{table}

% -----------------------------------------------------------------------------
\section{Phase~2: Deep Layer-A Semantic Fidelity Audit (\texorpdfstring{$n = 20$}{n = 20})}
\label{sec:phase2}

We ran Stage~5 adversarial Layer-A audits (\texttt{pro}-tier read-only referees) across all \textbf{20 locally inspected \texttt{engineering\_math\_proof} papers} in T2 (the 8 Pilot~v1 repos, \texttt{2611.15128} inside \texttt{leanVM}, 8 GitHub/Zenodo \texttt{math.OC} repos, and all 3 statically unpacked arXiv-ancillary/zipped packages \texttt{2611.26140}, \texttt{2611.16781}, and \texttt{2611.12585}; full theorem-by-theorem tables with \texttt{file:line} citations are archived in \texttt{reports/slices/2026-11/phase2\_layer\_a\_audit.md}). Figure~\ref{fig:phase2_heatmap} visualizes the complete severity matrix across Modes~\#1--\#4, non-vacuity witnesses, and Mathlib anchoring.

\begin{figure}[t]
\centering
\includegraphics[width=\textwidth]{fig_phase2_heatmap.pdf}
\caption{Phase~2 adversarial Layer-A semantic fidelity matrix across all 20 locally inspected T2 \texttt{engineering\_math\_proof} repositories (\texttt{pro}-tier read-only auditors). The bold divider contrasts the sole bulk automated pipeline (\texttt{2611.26272}, \emph{Paper2Lean}/\texttt{OptPaper-20}) against the 19 independent author-owned engineering repositories.}
\label{fig:phase2_heatmap}
\end{figure}

Four headline findings emerge from the expanded 20-paper Layer-A audit:

\begin{enumerate}
  \item \textbf{A natural control experiment --- Author-Owned Formalizations (19/19 clean on severe/moderate drift) vs.\ Bulk Automated Pipeline (\texttt{2611.26272}, \emph{Paper2Lean}):} Across all 19 independent author-owned repositories, our adversarial auditors found \textbf{zero} moderate/severe Mode~\#1 statement drift, \textbf{zero} undisclosed Mode~\#3 definition smuggling, and \textbf{zero} moderate Mode~\#4 junk/vacuity bugs---with \textbf{7/19} achieving \texttt{high\_faithful} (including \texttt{2611.09050} with \texttt{none} across all four modes, \texttt{2611.13584} covering all 88 numbered heavy-ball theorems, \texttt{2611.18455} constructing the McKean--Vlasov SDE from primitive Brownian motion, \texttt{2611.05243} matching every regret constant verbatim, \texttt{2611.16781} proving unrestricted AC uniqueness and mechanical crown LIL across 34.6k~LoC, and \texttt{P4} flagship \texttt{2611.13872} \emph{LongTails}). By stark contrast, the \textbf{sole \texttt{material\_undisclosed\_drift\_or\_risk} finding} in the entire 20-paper cohort was \textbf{\texttt{2611.26272} (\emph{Paper2Lean} / \texttt{OptPaper-20}, \texttt{v1}; repo \texttt{TTony2019/OptPaper} commit \texttt{1bde49b})}, a 151k-LoC bulk autoformalization of 20 optimization papers. Even on \texttt{Moudafi\_Oliny}---which \texttt{2611.26272}'s own Table~1 scores as \texttt{11/11 (100\%)} faithful---\texttt{thm\_1.lean:1500} adds an undisclosed global Lipschitz hypothesis \texttt{hB\_lipschitz} that is strictly stronger than the paper's quasi-cocoercivity definition \texttt{IsCocoerciveWRT} (\texttt{def\_5.lean:28-29}, quantified only over $y \in (A+B)^{-1}(0)$). Furthermore, despite the paper's abstract and introduction stating that ``all identified discrepancies were repaired in the released projects'' with ``no unresolved discrepancies,'' released files from partially-scored Table~1 projects still retain concrete discrepancies: \texttt{AbbasDynamical2015.kadek\_klee\_strong\_convergence} (\texttt{rem\_3.lean:33-57}) takes weak convergence \texttt{hweak} as a premise rather than deriving it; \texttt{Guler\_1992} (\texttt{thm\_3.lean:45-46}) adds \texttt{[FiniteDimensional $\mathbb{R}$ H]} to an infinite-dimensional Hilbert-space theorem (acknowledged in its Lean docstring, lines~42--44, but omitted from the paper text); and \texttt{Attouch\_Adly} (\texttt{thm\_9.lean:11-19}) proves only bare monotonicity (\texttt{AntitoneOn}) rather than an $o(1/\ln t)$ rate. This directly corroborates Bastounis et al.'s Type~(i) warning inside our own blind empirical audit.
  \item \textbf{Persistent computational/algorithmic scope boundary (18/20 Mode~\#2):} Across subfields, authors reliably formalize core analytic/algebraic theorems while leaving big-$O$ runtime complexity, Turing-machine reductions (\texttt{2611.14475}, \texttt{2611.05136}, \texttt{2611.12680}), iterative tolerances (\texttt{2611.10902}), external companion premises (\texttt{2611.12585} conditioning nonlinear KdV bounds on companion premise \texttt{P1} while proving six tails plus \texttt{prop318} unconditionally), or game-to-certificate reductions (\texttt{2611.26140} kernel-replaying all 15 surface-cut/continuum certificates by \texttt{decide} while candidly keeping topological \texttt{GuardsCorridor} reductions on paper; \texttt{2611.15128} omitting Theorem~1.1's asymptotic bound; \texttt{2611.23762} covering three companion helper lemmas) outside Lean.
  \item \textbf{Strong Mathlib anchoring (13 strong, 7 mixed) and improving non-vacuity hygiene (13/20):} Non-vacuity witness adoption rose markedly in \texttt{math.OC} and applied analysis (all 3 newly unpacked packages include concrete instantiations or explicit witness suites---\texttt{2611.26140} certificate replays, \texttt{2611.16781} \texttt{nominalGlobalExample}/\texttt{robustGlobalExample}, and \texttt{2611.12585} \texttt{CompanionFinal/Probes.lean}---bringing overall T2 coverage from 3/8 in Pilot~v1 to \textbf{13/20 = 65.0\%}).
  \item \textbf{Spec-locks still require independent spec auditing (\texttt{2611.04876} vs.\ \texttt{2611.13872}):} While \texttt{P4} \texttt{comparator} lock \texttt{2611.13872} (\emph{LongTails}) achieved \texttt{high\_faithful}, our auditor caught an undisclosed public-interface omission in \texttt{P4} repo \texttt{2611.04876} (\texttt{v1}, commit \texttt{7da3e11}: \texttt{version-1.0/lean/Compiler.lean:1936-1978} omits standalone exact-correctness equality from the comparator-checked \texttt{Compiler.main\_theorem}, even though exact correctness is proved internally in \texttt{Proofs/}). Comparator enforces proof-to-spec integrity, but independent spec auditing remains essential. Exact arXiv versions and pinned Git/Zenodo/ancillary provenance for all 23 audited repositories (239/239 resolved in-bounds \texttt{file:line} references) are archived in \texttt{artifacts/phase2\_provenance.json}.
\end{enumerate}

% -----------------------------------------------------------------------------
\section{Longitudinal Inflection Curve (Nov~2025 \texorpdfstring{$\rightarrow$}{->} May~2026 \texorpdfstring{$\rightarrow$}{->} Nov~2026)}
\label{sec:longitudinal}

To measure whether late-2026 engineering Lean adoption represents a steady multi-year baseline or a sharp 2026 AI-agentic inflection, we ran our exact eight-category pipeline on two matched historical 19-day windows---\textbf{T0} (\texttt{2025-11-01}--\texttt{2025-11-19}, $n=2{,}435$) and \textbf{T1} (\texttt{2026-05-01}--\texttt{2026-05-19}, $n=3{,}298$)---alongside our \textbf{T2} reference slice (\texttt{2026-11-01}--\texttt{2026-11-19}, $n=4{,}104$), totaling \textbf{9,837} unique preprints (\textbf{9,806} full-text screened = \textbf{99.68\%} coverage; full report in \texttt{reports/longitudinal\_analysis.md}).

\subsection{Eliminating Two Historical-arXiv Measurement Biases}
Constructing an unbiased historical baseline exposed two methodological traps that would severely distort naive longitudinal scraping:
\begin{enumerate}
  \item \textbf{OAI-PMH datestamp survivorship bias:} arXiv's OAI-PMH \texttt{from}/\texttt{until} parameters filter by latest revision \texttt{datestamp}, not \texttt{v1} submission date. Passing \texttt{until=2025-11-22} silently drops every Nov~2025 preprint later revised during Dec~2025--Nov~2026. Running a single-pass \texttt{arXivRaw} harvest with \texttt{from=2025-11-01} and \emph{unbounded \texttt{until}}---filtering strictly on \texttt{v1} \texttt{submitted\_date} and verifying against \texttt{arxiv.org/list}---rescued \textbf{862/2,435 (35.4\%)} revised T0 preprints and \textbf{1,048/3,298 (31.8\%)} revised T1 preprints.
  \item \textbf{Retroactive post-hoc autoformalization (\texttt{P1-retro}):} Screening latest versions (\texttt{v\_latest}) catches papers whose authors returned 5--7 months later to add Lean~4 formalizations once late-2026 AI proving agents matured. Direct \texttt{v1} HTML/PDF inspection of all revised historical Lean hits revealed \textbf{three} such cases with \emph{zero} Lean mentions in \texttt{v1}: \texttt{2511.15517} (\emph{Beluga}~\cite{arxiv2511_15517}, \texttt{cs.CR}/\texttt{cs.DC}; \texttt{v1} Nov~2025 $\rightarrow$ \texttt{v2} Jun~2026, Lean~\texttt{v4.33.0}), \texttt{2605.09552} (\emph{Phases of Muon}~\cite{arxiv2605_09552}, \texttt{math.OC}; \texttt{v1} May~2026 $\rightarrow$ \texttt{v2} Oct~2026, 340.3k~LoC Lean~\texttt{v4.34.1}), and \texttt{2605.19531} (\emph{Conflict-Freedom}~\cite{arxiv2605_19531}, \texttt{cs.DC}; \texttt{v1} May~2026 $\rightarrow$ \texttt{v2} Nov~2026, 34.3k~LoC Claude Opus~5.5-assisted Lean~\texttt{v4.34.0}). Without \texttt{v1} separation, naive latest-version scraping would misattribute \textbf{374.6k~LoC} of Oct--Nov~2026 code back to May~2026, inflating T1's paper-owned LoC by $105\times$. We therefore report both contemporaneous (\texttt{v1}) and latest-revision (\texttt{v\_latest}) rates.
\end{enumerate}

\subsection{The 2026 Inflection vs.\ Flat Beyond-Lean Formal Methods}
Table~\ref{tab:longitudinal} and Figure~\ref{fig:longitudinal} summarize the 12-month trajectory across all five comparative dimensions:
\begin{itemize}
  \item \textbf{Exponential Lean surge vs.\ flat traditional FV control:} Contemporaneous (\texttt{v1}) confirmed Lean adoption jumped from \textbf{0/2,435 (0.00\%}, one-sided 95\% Clopper--Pearson upper bound 0.12\%) in Nov~2025 to \textbf{5/3,298 (0.15\%}, 95\% Wilson CI $[0.06\%, 0.35\%]$) in May~2026 to \textbf{39/4,104 (0.95\%}, 95\% CI $[0.70\%, 1.30\%]$) in Nov~2026---\textbf{a $6.3\times$ jump in six months} ($23.1\times$ YoY on \texttt{v\_latest}: $0.04\% \rightarrow 0.21\% \rightarrow 0.95\%$). Over the exact same 9,837 papers, deterministic Stage-2 full-text hits for \emph{all other formal verification tools combined} (Coq/Rocq, Isabelle, Agda, Dafny, TLA+, UPPAAL, Tamarin, ProVerif, SMT) stayed flat-to-slightly-declining (\textbf{1.52\% $\rightarrow$ 1.33\% $\rightarrow$ 1.17\%}, $n=37 \rightarrow 44 \rightarrow 48$), while Lean~4 full-text candidates surged $10.9\times$ (\textbf{0.12\% $\rightarrow$ 0.58\% $\rightarrow$ 1.34\%}, $n=3 \rightarrow 19 \rightarrow 55$), \textbf{overtaking all other formal verification tools combined in Nov~2026}.
  \item \textbf{Disciplinary divergence:} \texttt{cs.CR} (0.44\%, 3/685~\cite{arxiv2605_02856,arxiv2605_03034,arxiv2605_16407}), \texttt{cs.DC} (0.32\%~\cite{arxiv2605_05801}), \texttt{eess.SP} (0.24\%~\cite{arxiv2605_14610}), and \texttt{eess.SY} (0.22\%~\cite{arxiv2605_03034}) hosted the earliest contemporaneous \texttt{v1} adopters in May~2026, whereas \textbf{\texttt{math.OC} ($0.00\% \rightarrow 0.00\%\,\texttt{v1}\,[0.20\%\,\texttt{v\_lat}] \rightarrow 2.62\%$)} and \textbf{\texttt{cs.IT} ($0.00\% \rightarrow 0.00\% \rightarrow 1.78\%$)} erupted exclusively in H2~2026. Across all three slices combined, \textbf{\texttt{cs.NI} (0/559)} and \textbf{\texttt{cs.RO} (0/2,431)} remain strictly at 0.00\%.
  \item \textbf{38.7$\times$ scale explosion, AI transition, and tier maturation:} Contemporaneous May~2026 Lean repositories were compact, human-scale artifacts (\textbf{median 749~LoC}, range 473--2,381~LoC, total 3,603~LoC across 3 repos), included the study's sole \texttt{explicitly\_no} human-authored Lean paper (\texttt{2605.14610}~\cite{arxiv2605_14610} in \texttt{eess.SP}, 473~LoC) and sole \texttt{P2\_lean\_first} paper (\texttt{2605.02856}~\cite{arxiv2605_02856} in \texttt{cs.CR}, 749~LoC, Claude Opus~4.7), had \textbf{40\%} explicit AI disclosure, and recorded \textbf{0/5 \texttt{Tier3\_Gold}}. By Nov~2026, median paper-owned LoC jumped \textbf{38.7$\times$} to \textbf{29,005~LoC} (total \textbf{3.10M} subpackage-adjusted paper-owned LoC; 4.30M raw repo LoC), explicit AI disclosure reached \textbf{76.9\%} (0\% human-only), \texttt{engineering\_math\_proof} rose from 3 (60.0\%) to 28 (71.8\%), and \texttt{Tier3\_Gold} reached \textbf{48.7\%} (19/39; notably, both \texttt{Tier3\_Gold} papers in T1 were the Oct--Nov~2026 retroactive \texttt{v2} revisions).
  \item \textbf{Layer-A semantic fidelity maturation (3 T1 audits, bringing total to $n=23$):} Auditing all three cloned T1 \texttt{engineering\_math\_proof} repositories (\texttt{pro}-tier referees, \texttt{reports/slices/2026-05/phase2\_layer\_a\_audit.md}; pinned SHAs in \texttt{artifacts/phase2\_provenance.json}) reveals a clear qualitative transition: both contemporaneous May~2026 \texttt{v1} Lean adopters (\texttt{2605.14610} [\texttt{v5}, commit \texttt{69788a9}] and \texttt{2605.02856} [\texttt{v2}, commit \texttt{915da15}]) scored \texttt{partial\_or\_weakened\_formalization} with \texttt{moderate} Mode~\#3 custom abstractions (e.g., \texttt{2605.14610} hardcodes tabular entropy formulas instead of Mathlib \texttt{MeasureTheory} integrals and proves $k_{\text{Lap}}=e/\sqrt{2}\approx 1.9221$ vs.\ the paper table's $1.9300$; \texttt{2605.02856} postulates a closed-form two-branch Montgomery map and omits induction in its $k$-stage composition). By contrast, T1's Oct~2026 retroactive \texttt{v2} (\texttt{2605.09552}, \emph{Phases of Muon}, commit \texttt{dc0ec6f}, 340.3k~LoC + 14.6k \texttt{alea-foundations}) achieved \texttt{mostly\_faithful\_with\_disclosed\_gaps} with \texttt{none} on both Mode~\#3 and Mode~\#4 (discharging all four external Gaussian inequalities and compact-group Haar measure directly from Mathlib, alongside explicit \texttt{phaseOmegaTimesOK} anti-junk guards)---matching the \textbf{19/19} independent Nov~2026 author-owned repos (\textbf{0/20} moderate/severe Mode~\#1 or Mode~\#3 failures across all Oct--Nov~2026 author-owned repos).
\end{itemize}

\begin{table}[t]
\centering
\caption{Longitudinal comparison across three matched 19-day eight-category arXiv slices ($N = 9{,}837$ unique preprints). Both contemporaneous (\texttt{v1}) and latest-revision (\texttt{v\_latest}, including retroactive 2026 \texttt{v2} Lean additions) figures are shown.}
\label{tab:longitudinal}
\scriptsize
\setlength{\tabcolsep}{3.4pt}
\begin{tabular}{lrrrr}
\toprule
\textbf{Metric} & \textbf{T0: Nov 2025} & \textbf{T1: May 2026} & \textbf{T2: Nov 2026} & \textbf{Shift (T1$\rightarrow$T2)} \\
\midrule
Corpus size (8 cats, unique) & 2,435 & 3,298 & 4,104 & $+24.4\%$ ($+68.5\%$ YoY) \\
Full-text coverage (HTML5+PDF) & 99.43\% & 99.70\% & 99.83\% & $\ge 99.4\%$ all slices \\
Later-revised preprints (\texttt{v>1}) & 862 (35.4\%) & 1,048 (31.8\%) & 248 (6.0\%) & Rescued (unbounded OAI) \\
\midrule
Stage-2 Beyond-Lean FV control & 37 (1.52\%) & 44 (1.33\%) & 48 (1.17\%) & Flat / slightly down \\
Stage-2 Lean candidates (full-text) & 3 (0.12\%) & 19 (0.58\%) & \textbf{55 (1.34\%)} & \textbf{2.3$\times$} (\textbf{10.9$\times$} YoY) \\
\textbf{Confirmed Lean (\texttt{v1} contemp.)} & \textbf{0 (0.00\%)} & \textbf{5 (0.15\%)} & \textbf{39 (0.95\%)} & \textbf{6.3$\times$ rate} \\
Confirmed Lean (\texttt{v\_latest}) & 1 (0.04\%) & 7 (0.21\%) & 39 (0.95\%) & 4.5$\times$ (23.1$\times$ YoY) \\
\texttt{engineering\_math\_proof} (\texttt{v1}/\texttt{v\_lat}) & 0 / 0 & 3 / 4 & \textbf{28} & \textbf{7.5$\times$} rate (\texttt{v1}) \\
\texttt{math.OC} Lean (\texttt{v1}/\texttt{v\_lat}) & 0.00\% & 0.00\% / 0.20\% & \textbf{2.62\% (19)} & Erupted H2~2026 \\
\midrule
Explicit AI disclosure (\texttt{v1} y/u/n) & 0 / 0 / 0 & 2 / 2 / 1 (40\%) & 30 / 9 / 0 (\textbf{76.9\%}) & Human-only $\rightarrow$ 0\% \\
Paradigms \texttt{v1} (\texttt{P1/P2/P3/P4/P0}) & 0/0/0/0/0 & 3/1/0/0/1 & 28/0/0/3/8 & \texttt{P4} emerges; \texttt{P3}=0 \\
Tiers \texttt{v1} (Gold/Silver/Bronze) & 0 / 0 / 0 & 0 / 2 / 3 & \textbf{19} / 7 / 13 & Gold: 0\% $\rightarrow$ 48.7\% \\
Median paper-owned LoC (\texttt{v1}) & \textemdash{} & 749 ($n=3$) & \textbf{29,005} ($n=28$) & \textbf{38.7$\times$} scale jump \\
Total paper-owned LoC (\texttt{v1}) & 0 & 3,603 & \textbf{3,104,750} & 3.6k $\rightarrow$ 3.10M \\
\bottomrule
\end{tabular}
\end{table}

\begin{figure}[t]
\centering
\includegraphics[width=\textwidth]{fig_longitudinal_curve.pdf}
\caption{Twelve-month longitudinal adoption curve across 9,837 engineering and applied-theory preprints (T0: Nov~1--19, 2025; T1: May~1--19, 2026; T2: Nov~1--19, 2026). (a) Contemporaneous (\texttt{v1}, with 95\% CIs) and latest-revision confirmed Lean~4 rates plus Stage-2 Lean candidates surge past the flat Beyond-Lean formal-methods control series. (b) Contemporaneous \texttt{v1} adoption trajectories across all eight categories, showing H2~2026 eruption in \texttt{math.OC} and \texttt{cs.IT} alongside persistent 0.00\% adoption in \texttt{cs.NI} and \texttt{cs.RO}.}
\label{fig:longitudinal}
\end{figure}

% -----------------------------------------------------------------------------
\section{Sandboxed Dynamic Kernel Verification (\texorpdfstring{$n = 14$}{n = 14} Compact Repositories)}
\label{sec:sandbox}

A natural reviewer objection to any static screening study is whether AI-generated engineering repositories actually compile in a fresh environment and whether their exported theorems are genuinely free of transitive \texttt{sorryAx}, \texttt{Lean.ofReduceBool} (\texttt{native\_decide}), or custom \texttt{axiom}s. To answer this directly without violating our static-analysis security boundary, we built an isolated two-phase Docker verification harness (\texttt{sandbox/}, image \texttt{lean-survey-sandbox:v1}) and evaluated the \textbf{exhaustive cohort of all 14 standalone primary Lean~4 repositories under 15,000 non-empty LoC} across T2 ($n=11$) and T1 ($n=3$, covering 100\% of contemporaneous May~2026 scanned repos; excluding only \texttt{2611.15128}, whose 4.96k-LoC subpackage lives inside the $>100\text{k}$-LoC \texttt{leanVM} monorepo).

\subsection{Two-Phase Clean-Room Security Architecture}
Untrusted Lean repositories can execute arbitrary code via \texttt{lakefile.lean}, \texttt{\#eval}, or Git hooks (Mode~\#7). Our harness enforces a strict separation between trusted dependency caching and isolated execution:
\begin{enumerate}
  \item \textbf{Phase~A (Trusted Toolchain \& Mathlib Cache Prep --- zero untrusted code mounted):} The host orchestrator reads only \texttt{lean-toolchain} (strictly regex-validated) and \texttt{lake-manifest.json} (pure JSON parsing; never invoking \texttt{lake} on the target repo), installs the pinned Lean toolchain via \texttt{elan}, fetches precompiled Mathlib \texttt{.olean} caches via \texttt{lake exe cache get} inside a synthetic empty project, and clones any extra manifest Git dependencies with \texttt{core.hooksPath=/dev/null}.
  \item \textbf{Phase~B (Isolated Execution \& Programmatic \texttt{\#print axioms}):} Each target repository runs serially inside an ephemeral container (\texttt{docker run --rm}) with \texttt{--network=none}, unprivileged \texttt{--user 65534:65534}, \texttt{--cap-drop=ALL}, \texttt{--security-opt=no-new-privileges:true}, \texttt{--cpus=4 --memory=10g --pids-limit=512}, and read-only mounts of both the source repo and cached dependencies. After \texttt{lake build} (or direct \texttt{lean} compilation for Std-only \texttt{2611.19551} per its \texttt{formal/README.md}) succeeds, \texttt{sandbox/check\_axioms.lean} executes Lean's kernel-level \texttt{Lean.collectAxioms} across \emph{every} non-internal theorem in every built project module, emitting structured JSON archived in \texttt{artifacts/sandbox\_builds/}.
\end{enumerate}

\begin{table}[t]
\centering
\caption{Isolated clean-room Docker \texttt{lake build} and programmatic \texttt{\#print axioms} (\texttt{Lean.collectAxioms}) results across all 14 standalone primary Lean~4 repositories $<15\text{k}$~LoC (T2: Nov~2026, $n=11$; T1: May~2026, $n=3$). All untrusted builds ran with \texttt{--network=none}, UID \texttt{65534}, \texttt{--cap-drop=ALL}, and read-only source mounts. \emph{Std only} = dependencies restricted to \{\texttt{propext}, \texttt{Classical.choice}, \texttt{Quot.sound}\}; \emph{synth.} = secondary \texttt{--network=none} check in the README/header-documented Lean+Mathlib release for repos omitting \texttt{lakefile}/\texttt{lean-toolchain}.}
\label{tab:sandbox}
\tiny
\setlength{\tabcolsep}{2.3pt}
\begin{tabularx}{\textwidth}{llllrllrrX}
\toprule
\textbf{arXiv ID} & \textbf{Sl.} & \textbf{Tier} & \textbf{Role} & \textbf{LoC} & \textbf{Toolchain} & \textbf{Ref} & \textbf{Time (s)} & \textbf{Thms} & \textbf{OOTB Outcome \& Transitive Axioms} \\
\midrule
2611.18220 & T2 & T3 & Eng & 528 & \emph{none} & \texttt{8e6dede} & \textemdash{} & 0 (27*) & \textbf{FAIL (packaging)}: no lakefile; synth.\ \texttt{v4.35.0} pass (80.4s, Std only) \\
2611.23762 & T2 & T2 & Eng & 589 & \emph{none} & \texttt{5e07f3a} & \textemdash{} & 0 (37*) & \textbf{FAIL (packaging)}: no lakefile; synth.\ \texttt{v4.36.0-rc2} pass (27.2s, Std only) \\
2611.14756 & T2 & T3 & Eng & 670 & \texttt{v4.35.0} & \texttt{5e25ee8} & 22.9 & 32 & \textbf{PASS} --- Std only (0 \texttt{sorryAx}, 0 \texttt{ofReduceBool}) \\
2611.19551 & T2 & T3 & Sys & 1,332 & \texttt{v4.19.0} & \texttt{f2121df} & 15.8 & 68 & \textbf{PASS} (Std-only \texttt{lean}) --- Std only (7 formal files) \\
2611.10902 & T2 & T3 & Eng & 2,068 & \texttt{v4.36.0-rc2} & \texttt{fb87e63} & 60.3 & 113 & \textbf{PASS} --- Std only \\
2611.02256 & T2 & T2 & TCS & 2,191 & \texttt{v4.36.0-rc1} & \texttt{1f9e89b} & 21.0 & 193 & \textbf{PASS} --- Std only \\
2611.15201 & T2 & T3 & TCS & 5,307 & \texttt{v4.35.0-rc2} & \texttt{z/24127514} & 43.8 & 274 & \textbf{PASS} (\texttt{SparseMM}+\texttt{Solution}) --- Std only \\
2611.16690 & T2 & T2 & TCS & 6,592 & \texttt{v4.35.0-rc3} & \texttt{2318a73} & 59.8 & 570 & \textbf{PASS (Mode~\#6)}: \texttt{sorryAx} in 9 thms (incl.\ \texttt{theorem\_3\_2}) \\
2611.10810 & T2 & T3 & Sys & 10,862 & \texttt{v4.32.2} & \texttt{z/24035424} & 11.2 & 507 & \textbf{PASS} (zero external deps) --- Std only \\
2611.05243 & T2 & T3 & Eng & 11,204 & \texttt{v4.33.1} & \texttt{49ade84} & 129.7 & 715 & \textbf{PASS} --- Std only \\
2611.18455 & T2 & T3 & Eng & 12,236 & \texttt{v4.34.1} & \texttt{e8efa5d} & 1,562.1 & 554 & \textbf{PASS} (builds \texttt{stochasticAnalysis}) --- Std only \\
\midrule
2605.14610 & T1 & T2 & Eng & 473 & \texttt{v4.26.0} & \texttt{69788a9} & 28.4 & 25 & \textbf{PASS} --- Std only \\
2605.02856 & T1 & T2 & Eng & 749 & \texttt{v4.30.0-rc1} & \texttt{915da15} & 0.9 & 0 & \textbf{FAIL (local dep)}: missing \texttt{../qanary-paper6} (+3 siblings) \\
2605.16407 & T1 & T1 & Sys & 2,381 & \texttt{v4.30.0-rc2} & \texttt{8e5b718} & 24.5 & 139 & \textbf{PASS (Mode~\#6)}: 11 custom \texttt{axiom}s (incl.\ Cauchy--Schwarz, Hoeffding) \\
\bottomrule
\multicolumn{10}{@{}p{\textwidth}@{}}{\footnotesize *Secondary synthetic-environment check using the exact Lean+Mathlib release stated in the repo README/header (total audited: 3,190 OOTB; 3,254 incl.\ synth.).}
\end{tabularx}
\end{table}

\begin{figure}[t]
\centering
\includegraphics[width=\textwidth]{fig_sandbox_builds.pdf}
\caption{Sandboxed dynamic kernel verification across all 14 standalone compact Lean~4 repositories ($<15\text{k}$~LoC across T1 and T2; \texttt{--network=none}, unprivileged container). (a) Per-repository container build time (log scale), audited theorem counts, and transitive axiom outcomes. (b) Breakdown of dynamic ground-truth verdicts by static verification tier, confirming 7/7 axiom-clean out-of-the-box \texttt{Tier3\_Gold} builds and 100\% agreement with static Mode~\#6 downgrades.}
\label{fig:sandbox}
\end{figure}

\subsection{Dynamic Benchmark Findings}
Table~\ref{tab:sandbox} and Figure~\ref{fig:sandbox} establish four empirical results that directly ground-truth our static screening pipeline:
\begin{enumerate}
  \item \textbf{High out-of-the-box reproducibility (78.6\%) with zero proof-compilation errors:} Across all \textbf{14} compact repositories, \textbf{11/14 (78.6\%)} build cleanly out of the box (\textbf{9/11 = 81.8\%} in T2; \textbf{2/3 = 66.7\%} in T1), auditing \textbf{3,190} distinct theorems (median build wall time \textbf{28.4\,s}). Crucially, \textbf{0/14} suffer genuine Lean proof or type-checking errors: all three out-of-the-box failures stem from \emph{packaging hygiene} (Mode~\#12). Specifically, \texttt{2611.18220} (\texttt{Lean-SysID}) and \texttt{2611.23762} (\texttt{CoprimeFrobenius}) omit \texttt{lakefile} and \texttt{lean-toolchain}; when placed inside a \texttt{--network=none} synthetic project matching their documented Lean+Mathlib versions (\texttt{v4.35.0} and \texttt{v4.36.0-rc2}), \textbf{2/2} compile cleanly with exclusively standard axioms (\textbf{27} and \textbf{37} theorems, bringing total audited theorems to \textbf{3,254}). The third failure, T1 repo \texttt{2605.02856} (\emph{Qanary}), aborts in 0.86\,s because \texttt{lakefile.lean} hardcodes four unbundled relative sibling paths (\texttt{../qanary-paper6}, etc.).
  \item \textbf{Static \texttt{Tier3\_Gold} labels accurately predict clean kernel verification:} Among all \textbf{8} compact T2 \texttt{Tier3\_Gold} repositories, \textbf{7/8 (87.5\%)} build cleanly out of the box and \textbf{7/7 (100\%)} pass programmatic \texttt{Lean.collectAxioms} across \textbf{2,263} theorems with zero transitive \texttt{sorryAx}, zero \texttt{Lean.ofReduceBool}, and zero custom axioms. The sole \texttt{Tier3\_Gold} out-of-the-box failure (\texttt{2611.18220}) was already flagged with static \texttt{Mode12\_ProcessOpacity} and passes axiom-clean under its README's \texttt{v4.35.0} toolchain.
  \item \textbf{Exact dynamic corroboration of Mode~\#6 and zero \texttt{Lean.ofReduceBool}:} Exactly \textbf{2/11} out-of-the-box passing repos exhibit nonstandard axioms, and \emph{both} were already downgraded below \texttt{Tier3\_Gold} with static Mode~\#6 flags: \texttt{2611.16690} (\texttt{Tier2\_Silver}) compiles cleanly across \textbf{570} theorems but propagates \texttt{sorryAx} into 9 theorems (including \texttt{AnyDim.Paper.theorem\_3\_2}, \texttt{corollary\_3\_3}, and \texttt{corollary\_3\_4}), while T1 repo \texttt{2605.16407} (\texttt{Tier1\_Bronze}, \textbf{139} theorems) postulates \textbf{11} unproved custom axioms---including standard mathematical inequalities (\texttt{cauchy\_schwarz\_sq}, \texttt{HoeffdingInequality}) alongside oracle assumptions. Across all 14 repositories, \textbf{0/14} depend on \texttt{Lean.ofReduceBool} or \texttt{Lean.trustCompiler}.
  \item \textbf{Extreme toolchain fragmentation and uncached-dependency latency:} Remarkably, all \textbf{12} pinned repositories require \textbf{12 distinct Lean~4 toolchains} spanning \texttt{v4.19.0} through \texttt{v4.36.0-rc2} (including 5 release candidates), underscoring why automated clean-room reproduction requires multi-toolchain caching. While Mathlib-cached and zero-dependency repos build rapidly (\textbf{11.2\,s} for 10.9k-LoC zero-dep \texttt{2611.10810} to 129.7\,s for 11.2k-LoC \texttt{2611.05243}), uncached third-party Git dependencies dominate latency: \texttt{2611.18455} (\texttt{wmfki\_lean}) required \textbf{1,562.1\,s} (26.0\,min) primarily to compile its external \texttt{stochasticAnalysis} dependency from source inside the offline sandbox.
\end{enumerate}

% -----------------------------------------------------------------------------
\section{Threats to Validity and Next Steps}
\label{sec:limitations}

Three primary limitations remain and define our next engineering milestones:
\begin{itemize}
  \item \textbf{Dynamic verification scale ceiling \& human semantic sign-off:} While Section~\ref{sec:sandbox} dynamically verifies all \textbf{14} compact standalone repositories ($<15\text{k}$~LoC) in isolated containers, larger multi-10k to 929k-LoC codebases remain statically audited due to multi-hour build budgets (and we did not yet execute \texttt{lake comparator} or external multi-kernel replay). Likewise, Phase~2 Layer-A LLM semantic auditing ($n=23$) serves as a high-leverage referee pre-filter rather than a substitute for independent human mathematician sign-off.
  \item \textbf{Remaining unreachable external hosts:} While our 9,837-paper longitudinal design resolves the single-snapshot limitation, Stage-2 recall is calibrated ($\le 1.49\%$ 95\% upper bound on GitHub-linked negatives), and all three arXiv ancillary/zipped packages (\texttt{2611.26140}, \texttt{2611.16781}, \texttt{2611.12585}) are now both statically profiled (all zero \texttt{sorry}/custom \texttt{axiom}) and Phase-2 Layer-A audited, one institutional GitLab host (\texttt{2611.22732}, \texttt{gitlab.mpi-sws.org}) persistently timed out on port~443 upon retry (alongside external portal \texttt{2611.09169}).
  \item \textbf{ML theory scope:} \texttt{cs.LG} and \texttt{stat.ML} only entered this study when cross-listed into \texttt{math.OC}/\texttt{cs.CR} (already yielding two \texttt{high\_faithful}, dynamically verified \texttt{Tier3\_Gold} Lean papers, \texttt{2611.05243} and \texttt{2611.18455}); a dedicated ML-theory arm remains to be measured.
\end{itemize}

\paragraph{Next Milestones:} (1) \textbf{Human inter-rater reliability (IRR)} calibration on Stage-4 classifications and human sign-off on Stage-5 Layer-A named findings plus clean-cohort spot-checks; (2) extending sandboxed Docker verification to \textbf{medium/large ($>15\text{k}$~LoC) repos and \texttt{lake comparator} replay}; (3) a \textbf{stratified \texttt{cs.LG}/\texttt{stat.ML} theory arm} to benchmark ML theory against \texttt{math.OC} and \texttt{cs.IT}; and (4) a \textbf{standalone periodic runner} (\texttt{agy -p}) for unattended quarterly longitudinal snapshots.

\section*{AI-Use Disclosure and Human Accountability}
To uphold transparency standards for empirical LLM-assisted research, we document the exact harness, model, and reasoning-effort configuration verified from execution transcripts across all tasks (November~21--23, 2026), separating AI used as an empirical research instrument from AI used in software engineering and manuscript preparation:

\begin{table}[!htbp]
\centering
\caption{Exact AI harness, model, and reasoning-effort provenance verified from execution transcripts across all empirical and authoring tasks (November~21--23, 2026).}
\label{tab:ai_provenance}
\scriptsize
\setlength{\tabcolsep}{3.0pt}
\begin{tabularx}{\textwidth}{@{}>{\raggedright\arraybackslash}p{3.1cm}>{\raggedright\arraybackslash}p{2.85cm}>{\raggedright\arraybackslash}p{4.55cm}>{\raggedright\arraybackslash}X@{}}
\toprule
\textbf{Task / Pipeline Role} & \textbf{Harness \& Alias} & \textbf{Exact Model \& Effort} & \textbf{Scope \& Archived Provenance} \\
\midrule
\textbf{Stage 4: Lean Classification} (Pilot~v1/v2, T0, T1) & Google Antigravity \texttt{run\_workflow} (\texttt{flash}) & \textbf{Gemini 4 Hydrogen (High)}\newline \texttt{gemini-4-hydrogen-tiered-high} & 80 read-only subagent runs (30 Pilot~v1 + 50 Pilot~v2/T0/T1); \texttt{workflows/classify\_candidates.py}. \\
\textbf{Stage 4: Beyond-Lean Baseline} (Pilot~v1) & Google Antigravity \texttt{run\_workflow} (\texttt{flash\_lite}) & \textbf{Gemini 3.5 Flash-Lite}\newline \texttt{gemini-3.5-flash-lite} & 34 read-only subagent runs across the 6 Pilot~v1 categories; \texttt{workflows/classify\_candidates.py}. \\
\textbf{Stage 5: Layer-A Semantic Audit} (T2 + T1) & Google Antigravity \texttt{run\_workflow} (\texttt{pro}) & \textbf{Gemini 4 Hydrogen (Extra High)}\newline \texttt{gemini-4-hydrogen-tiered-xhigh} & 23 read-only adversarial audits (8 Pilot~v1 + 9 Pilot~v2 + 3 ancillary/zip + 3 T1); \texttt{workflows/audit\_layer\_a.py}. \\
\textbf{Code Writing \& Manuscript Drafting} (All Sections) & Google Antigravity CLI harness & \textbf{Gemini 4 Neon (High)}$^*$\newline \texttt{gemini-4-neon-high} & All pipeline, sandbox, analysis, audit, and figure code; drafting partner for Abstract and Sections~1--8. \\
\bottomrule
\multicolumn{4}{@{}l}{\footnotesize $^*$Auxiliary interactive repository sessions also used Gemini~4~Neon (Medium) and Gemini~4~Carbon (High).}
\end{tabularx}
\end{table}

\begin{enumerate}
  \item \textbf{AI as an empirical research instrument (Stages~4--5):} Deterministic zero-token scripts execute all ingestion, full-text screening, static artifact profiling, Stage-2 recall calibration, and isolated Docker kernel builds (Stages~1--3 and Section~\ref{sec:sandbox}). Stages~4--5 dispatched strictly read-only, non-executing \texttt{research} agents under strict JSON schemas via Google Antigravity's \texttt{run\_workflow} orchestrator (Table~\ref{tab:ai_provenance}); full prompt templates, schemas, and structured JSON outputs are archived verbatim in \texttt{workflows/} and \texttt{artifacts/}.
  \item \textbf{AI-assisted software engineering and manuscript preparation:} The Google Antigravity harness running \textbf{Gemini~4~Neon (High effort)} served as the coding engine for all pipeline, sandbox, analysis, audit, and figure-generation code and as a drafting partner across all text sections of this manuscript (Abstract and Sections~1--8). All empirical counts, code-generated vector figures, sandboxed kernel logs (\texttt{artifacts/sandbox\_builds/}), named repository critiques, pinned commit SHAs (\texttt{artifacts/phase2\_provenance.json}), live arXiv bibliography records (\texttt{scripts/audit\_bib.py --live}), and final manuscript claims were reviewed, verified, and edited by the human authors, who take full scholarly responsibility for the accuracy and integrity of this work.
\end{enumerate}

\bibliographystyle{plainnat}
\bibliography{../shared/references}

\end{document}
